% =====================================================================
%  RESUMEN EJECUTIVO EN INGLÉS  ---  1 PÁGINA / 200 A 400 PALABRAS
% =====================================================================
%  La rúbrica (categoría A, 2 puntos) exige CINCO criterios. Si falta uno solo,
%  bajas a "competente"; si faltan tres, quedas "con necesidad de mejora".
%
%  LOS CINCO CRITERIOS, EN ORDEN, Y CÓMO CUBRIRLOS:
%
%  [1] CONTEXTUALIZA EL LUGAR DE PRÁCTICA -> empresa, área y puesto de trabajo.
%      Los tres, explícitos. No basta con nombrar la empresa.
%        DTS SpA · R&D department (I+D+i) · Electrical Engineering intern.
%
%  [2] DESCRIBE EL OBJETIVO O INTERÉS DE LA EMPRESA a lograr al finalizar la
%      práctica. OJO: es el objetivo DE LA EMPRESA, no el del proyecto.
%        -> take the company's first formal step into space capabilities and
%           bring the payload design up to PDR maturity before committing
%           resources to detailed design.
%      Conviene declarar TAMBIÉN el objetivo del proyecto, en una oración
%      separada, para que ambos queden explícitos.
%
%  [3] INFORMACIÓN MÁS RELEVANTE DEL DOCUMENTO:
%        actividades principales -> CONOPS, requirements engineering and
%          traceability matrix, MBSE migration, orbital and link-budget simulator;
%        metodologías -> sequential systems-engineering life cycle with formal
%          design reviews (SRR, PDR) and the ARCADIA method on Eclipse Capella;
%        principales resultados -> 13 user requirements and 23 system
%          requirements traced, 8 high-risk items mitigated, 19 verification
%          tests specified, link closure on 4-5 daily passes with ~10 min
%          windows above a 14 dB Eb/N0 threshold, ROM estimate of USD 133,000;
%        conclusiones -> the architecture is feasible; SRR and PDR closed on time.
%
%  [4] CORRECTO USO DEL INGLÉS. Voz activa, tiempo pasado, sin traducción
%      literal del español. Revisa concordancia y ortografía.
%
%  [5] SÍNTESIS CLARA de contexto, objetivos, metodologías, resultados y
%      conclusiones. El corrector debe entender el informe completo leyendo
%      solo esta plana.
%
%  ESTRUCTURA SUGERIDA -- 5 párrafos, ~65 palabras cada uno (total ~330):
%    P1  Where and as what: company, area, role, dates. One sentence on what DTS
%        does and why the space line was new.
%    P2  The company's objective for the internship + the project's objective.
%    P3  Methodology: life-cycle reviews (IEEE Std 15288.2), ARCADIA/Capella,
%        and the custom Python link-budget simulator.
%    P4  Results, with numbers. This is the paragraph the corrector scans.
%    P5  Conclusion: feasibility demonstrated, PDR closed, what remains for CDR.
%
%  EJEMPLO DE ARRANQUE DE P1 (adáptalo, no lo copies literal):
%    "This report describes the professional internship carried out between
%     January 2 and February 28, 2026, at Desarrollo de Tecnologías y Sistemas
%     SpA (DTS), a Chilean electronics company serving the defense, security and
%     aerospace markets. The work took place in the R&D department as an
%     Electrical Engineering intern, within a space development line that the
%     company had not previously pursued."
%
%  EJEMPLO DE ARRANQUE DE P4:
%    "The internship produced a closed requirements baseline of 13 user
%     requirements and 23 system requirements, fully traced in both directions,
%     together with mitigation measures for the 8 items rated as high risk and a
%     preliminary verification matrix of 19 tests. The link-budget simulation
%     confirmed that the communication link closes on four to five daily passes
%     over Chilean territory, with contact windows of approximately ten minutes."
%
%  RECORDATORIO: escribe el abstract AL FINAL, cuando el cuerpo esté cerrado.
% =====================================================================

\selectlanguage{British}

\noindent{\Large\textbf{Abstract}\par\vspace{1ex}}

% >>> REDACTADO (389 palabras, dentro del rango 200-400 de la rúbrica).
%     Revisa el inglés y ajusta lo que no suene a como escribes tú.

This report describes the professional internship carried out between 2 January
and 28 February 2026 at Desarrollo de Tecnologías y Sistemas SpA (DTS), a
Chilean electronics company serving the defence, security and aerospace markets.
The work took place in its Research and Development department as an Electrical
Engineering intern, on a space development line defined in the company's
strategic plan but never executed.

DTS pursued two objectives during the period: to take the first formal step
towards developing space capabilities, and to make use of the technology
transfer agreement signed with INVAP. The constraint was not technical but one
of capacity: with no staff available to open a new line of work, the company
started it with one staff engineer and one intern. Both objectives converged on a single
deliverable: bringing the payload design up to Preliminary Design Review
maturity before committing resources to detailed design and manufacturing.

The project, named BIFROST, designs a Store \& Forward communications payload
for a 12U CubeSat that receives autonomous messages from remote ground stations
in UHF, stores them on board and downloads them over Chilean territory. The work followed a sequential systems engineering life cycle
structured around formal design reviews, with IEEE Std 15288.2-2014 defining the
deliverables and exit criteria of each gate. The system was modelled
with the ARCADIA method on Eclipse Capella, and a custom orbital and link budget
simulator was developed in Python.

The internship closed a baseline of 13 user and 21 system requirements linked by
35 traceability relations with full bidirectional coverage. Twelve were rated
high risk and mitigated, and 19 verification tests were specified. The link budget simulation confirmed
that the link closes above the 14 dB Eb/N0 threshold on four to five daily
passes, with contact windows of about ten minutes. It also revealed a zenith
null in the ground antenna pattern that constrains download planning, which a
constant-gain calculation would have missed.

Both the System Requirements Review and the Preliminary Design Review were
closed within the period, confirming the architecture as technically feasible. DTS retained a reusable system
model, a simulation tool and a documented methodology, began training its engineers in space
projects, and continues developing BIFROST towards the Critical Design Review.

\selectlanguage{Spanish}
\newpage
